\begin{figure}[ht]
\centering
\begin{tikzpicture}[x=1cm, y=1cm, font=\small]

% Axes
\draw[->, thick] (-3.6, 0) -- (3.6, 0) node[right, font=\scriptsize] {$x_{1}$};
\draw[->, thick] (0, -1.5) -- (0, 1.5) node[above, font=\scriptsize] {$x_{2}$};

% Elliptic contours of f(x_1, x_2) = 1/2 (x_1^2 + 100 x_2^2)
% level set f = c: x_1^2/(2c) + x_2^2/(c/50) = 1
% semi-axes (sqrt(2c), sqrt(c/50)). Choose c values that give visible ellipses.
\foreach \c in {0.5, 1.5, 4.5, 9, 13.5} {
  \pgfmathsetmacro{\axA}{sqrt(2*\c)}
  \pgfmathsetmacro{\axB}{sqrt(\c/50)}
  \draw[engblue, thin] (0,0) ellipse ({\axA} and {\axB});
}

% Gradient-descent zigzag trajectory from (3.0, 1.0) with eta = 0.018
% chosen so the iterates oscillate visibly across the long axis.
% Hand-computed iterates of (x_{k+1}) = (1 - eta) x_1, (1 - 100 eta) x_2
% with eta = 0.018: (1 - eta) = 0.982, (1 - 100 eta) = -0.8
\fill[engred] (3.00, 1.00) circle (1.6pt);
\fill[engred] (2.95, -0.80) circle (1.6pt);
\fill[engred] (2.89, 0.64) circle (1.6pt);
\fill[engred] (2.84, -0.51) circle (1.6pt);
\fill[engred] (2.79, 0.41) circle (1.6pt);
\fill[engred] (2.74, -0.33) circle (1.6pt);
\fill[engred] (2.69, 0.26) circle (1.6pt);
\fill[engred] (2.64, -0.21) circle (1.6pt);

\draw[engred, thick]
  (3.00, 1.00) -- (2.95, -0.80) -- (2.89, 0.64) -- (2.84, -0.51)
  -- (2.79, 0.41) -- (2.74, -0.33) -- (2.69, 0.26) -- (2.64, -0.21);

% Annotate starting and current points
\node[engred, anchor=west, font=\scriptsize] at (3.05, 1.10) {$\vect{x}_{0}$};
\node[engred, anchor=west, font=\scriptsize] at (2.65, -0.30) {$\vect{x}_{7}$};

% Annotate the long and short axes
\node[anchor=west, font=\scriptsize] at (3.5, 0.05) {long axis};
\node[anchor=west, font=\scriptsize] at (0.15, 1.2)
  {short axis (steep curvature)};

% Title under the figure body
\node[anchor=north, font=\scriptsize, align=center] at (0, -1.9)
  {$f(x_{1}, x_{2}) = \tfrac{1}{2}(x_{1}^{2} + 100 x_{2}^{2})$,
   $\kappa = 100$};

\end{tikzpicture}
\caption[Gradient-descent zigzag on an elongated quadratic]{Gradient
descent with a fixed step size $\eta$ on the elongated quadratic
$f(x_{1}, x_{2}) = \tfrac{1}{2}(x_{1}^{2} + 100 x_{2}^{2})$, plotted
on its elliptic level sets. The step size large enough to make
useful progress along the long axis ($x_{1}$, eigenvalue $1$)
overshoots the short axis ($x_{2}$, eigenvalue $100$) and the
iteration zigzags. The iteration count to reach a given accuracy
scales with the condition number $\kappa = L/\mu = 100$, the ratio
of the largest to the smallest Hessian eigenvalue. Newton's method
on the same problem reaches the optimum in one step; preconditioning
the gradient by an approximate inverse Hessian achieves the same
effect for general smooth objectives.}
\label{fig:vol02:ch15:gradient-zigzag}
\end{figure}
